\documentclass[10pt,a4paper]{amsart}
\usepackage[margin=19mm]{geometry}
\usepackage{amsmath,amssymb,mathtools}
\usepackage[scr=rsfs,cal=euler]{mathalfa}
\usepackage{xfrac}
\usepackage[unicode,hidelinks,bookmarks=false]{hyperref}
\newcommand{\ie}{i.e.\ }
\newcommand{\cf}{cf.\ }
\newcommand{\resp}{respectively\ }
\newcommand{\Subsection}{\subsection}
\providecommand{\nobreakdash}{\nobreak\hbox{-}}
\setlength{\emergencystretch}{2em}
\multlinegap0pt
\usepackage{amssymb}
\usepackage{mathrsfs}

\DeclareMathOperator{\SL}{\mathsf{SL}}
\DeclareMathOperator{\PSL}{\mathsf{PSL}}
\DeclareMathOperator{\Gr}{\mathcal{G}}
\DeclareMathOperator{\GL}{\mathsf{GL}}
\DeclareMathOperator{\SO}{\mathsf{SO}}
\DeclareMathOperator{\Sp}{\mathsf{Sp}}
\DeclareMathOperator{\idd}{id}
\DeclareMathOperator{\tr}{tr}
\DeclareMathOperator{\spann}{span}
\DeclareMathOperator{\dist}{dist}
\DeclareMathOperator{\rSL}{SL}
\DeclareMathOperator{\Ad}{Ad}
\DeclareMathOperator{\Aut}{Aut}


\newcommand{\clA}{\mathcal{A}}
\newcommand{\F}{\mathcal{F}}
\newcommand{\N}{\mathbb{N}}
\newcommand{\R}{\mathbb{R}}
\newcommand{\RR}{\R}
\newcommand{\ZZ}{\mathbb{Z}}
\newcommand{\C}{\mathbb{C}}
\newcommand{\HH}{\mathbb{H}}
\newcommand{\sG}{\mathsf{G}}
\newcommand{\scU}{\mathscr{U}}

\newcommand{\Es}{E^s}
\newcommand{\Ec}{E^c}
\newcommand{\Eu}{E^u}
\newcommand{\Ecs}{E^{cs}}
\newcommand{\Ecu}{E^{cu}}
\newcommand{\eps}{\varepsilon}
\newcommand{\RP}{{\RR\mathbb{P}}}
\renewcommand{\a}{\mathfrak{a}}
\newcommand{\ap}{\mathfrak{a}^+}
\newcommand{\liesl}{\mathfrak{sl}_2(\RR)}

\renewcommand*{\to}{\mathchoice{\longrightarrow}{\rightarrow}{\rightarrow}{\rightarrow}}

\let\originalleft\left 
\let\originalright\right
\renewcommand{\left}{\mathopen{}\mathclose\bgroup\originalleft}
\renewcommand{\right}{\aftergroup\egroup\originalright}
%%%%%%%%%%%%
\let\oldmapsto\mapsto
\renewcommand*{\mapsto}{\mathchoice{\longmapsto}{\oldmapsto}{\oldmapsto}{\oldmapsto}}

\let\oldin\in
\renewcommand{\in}{\mathchoice{\oldin}{\oldin}{\,\oldin\,}{\,\oldin\,}}

\newcommand*{\mk}{\mkern -1mu}
\newcommand*{\Mk}{\mkern -2mu}
\newcommand*{\mK}{\mkern 1mu}
\newcommand*{\MK}{\mkern 2mu}

\hypersetup{urlcolor=purple, linkcolor=blue, citecolor=red}

\newcommand*{\relabel}{\renewcommand{\labelenumi}{(\theenumi)}}
\newcommand*{\romanenumi}{\renewcommand*{\theenumi}{\roman{enumi}}\relabel}
\newcommand*{\Romanenumi}{\renewcommand*{\theenumi}{\Roman{enumi}}\relabel}
\newcommand*{\alphenumi}{\renewcommand*{\theenumi}{\alph{enumi}}\relabel}
\newcommand*{\Alphenumi}{\renewcommand*{\theenumi}{\Alph{enumi}}\relabel}
\let\oldtilde\tilde
\renewcommand*{\tilde}[1]{\mathchoice{\widetilde{#1}}{\widetilde{#1}}{\oldtilde{#1}}{\oldtilde{#1}}}
\let\oldhat\hat
\renewcommand*{\hat}[1]{\mathchoice{\widehat{#1}}{\widehat{#1}}{\oldhat{#1}}{\oldhat{#1}}}

\numberwithin{equation}{section}
\newtheorem{theo}{Theorem}[section]
\newtheorem{prop}[theo]{Proposition}
\newtheorem{coro}[theo]{Corollary}
\newtheorem{lemm}[theo]{Lemma}
\theoremstyle{definition}\newtheorem{defi}[theo]{Definition}
\newtheorem{rema}[theo]{Remark}
\title[A criterion for non-robust quasi-isometry]{Quasi-isometric free group representations into $\SL_3(\R)$: original Proposition 3.4}
\author{Le\'on Carvajales and Pablo Lessa and Rafael Potrie}
\date{Source: Annales Henri Lebesgue 8 (2025), 965--994; DOI 10.5802/ahl.252}
\begin{document}\maketitle
\noindent\textit{Editorial scope.} The counted unit is original Proposition 3.4 with its complete authored proof. The original eigenvalue-moduli and loxodromic notation of the same subsection, the recall of unipotent elements and the complete original Lemma 3.3 (openness of the power map) used in the proof are retained as uncounted context in original file order. The preceding generalities and the later applications to Barbot representations are omitted. Original published section, subsection, statement and display numbers are reproduced by the retained counters. Printed slips are kept literal without repair.
\setcounter{section}{2}
\section{Perturbed not quasi-isometric representations}\label{sec: generalities QI}

The most important property that we will need is the following. Recall that an element $g\neq 1$ of $\GL_d(\R)$ is said to be \emph{unipotent} if all its eigenvalues are equal to $1$.

\setcounter{subsection}{1}
\subsection{A criterion for non quasi-isometry}\label{ss.criterionnonQI}

For a matrix $g\in\SL_3(\R)$ we let
\[
\lambda_u(g)\geq\lambda_c(g)\geq\lambda_s(g)
\]
be the moduli of
the eigenvalues of $g$. Recall that $g$ is said to be \emph{loxodromic} if the inequalities above are all strict. In this case, the eigenvalues of $g$ are real, and denoted respectively by $\mu_u(g),\mu_c(g)$ and $\mu_s(g)$. The corresponding eigenspaces are one dimensional and denoted by $\Eu(g),\Ec(g)$
and $\Es(g)$. We also let
\[
\Ecu(g):=\Eu(g)\oplus\Ec(g) \quad\text{and}\quad
\Ecs(g):=\Es(g)\oplus\Ec(g),
\]
which are respectively the
attracting and repelling hyperplane of $g$ acting on the Grassmannian $\Gr_2(\R^3)$.

\setcounter{theo}{2}
\begin{lemm}\label{lem: power map is open}
Let $g_0\in\SL_3(\RR)$ be a loxodromic element and $n$ be a non-zero integer.
Then the map
\[
\SL_3(\RR)\to \SL_3(\RR): g\mapsto g^n
\]
sends every small enough neighborhood $U$ of $g_0$ to a neighborhood $U^{(n)}$ of $g_0^n$.
\end{lemm}


\begin{proof}
Indeed, every power of a loxodromic element is again loxodromic. Moreover, a neighborhood of a loxodromic element is parametrized by the choice of three lines $\Eu,\Ec$ and $\Es$ in general position, plus the choice of three non zero real numbers $\mu_u,\mu_c$ and $\mu_s$ of different moduli, and whose product is equal to one. As the eigenlines of a power of $g$ coincide with
eigenlines of $g$, and the map
\[
\R\setminus\{0\}\to \R\setminus\{0\}:
\mu\mapsto \mu^n
\]
is open, the proof is complete.
\end{proof}

\begin{prop}\label{prop: criterion for non robust qi}
Let $\Gamma$ be a non-abelian free group and $\rho:\Gamma\to \SL_3(\R)$ be a representation. Suppose that $\Gamma$ admits a free generating set containing two elements $a$ and $b$ with the following properties:
\begin{enumerate}
\item there exist non zero integers $m$ and $n$ so that the word $\omega:=a^mb^n$ satisfies that for some non zero integer $q$ and some $\mu\neq \pm 1$ the subspace $P_0:=\ker(\rho(\omega^q)-\mu)$ is two-dimensional. Moreover, the subspace $L_0:=\ker(\rho(\omega^q)-\mu^{-2})$ is one-dimensional.
\item The matrices $\rho(a)$ and $\rho(b)$ are loxodromic and $L_0$ is contained either in the attracting plane $\Ecu(\rho(a))$ of $\rho(a)$, or in the repelling one $\Ecs(\rho(a))$.
\end{enumerate}
Then there exists a continuous path $\rho_t:\Gamma\to\SL_3(\R)$ with $\rho_0=\rho$ and a sequence $t_k\to 0$ such that for every $k$ the element
\[
\rho_{t_k}([\omega^q,a^{p_k}\omega^qa^{-p_k}])
\]
is unipotent for some non zero integer $p_k$.
\end{prop}

\begin{proof}
Without loss of generality we assume that $L_0$ belongs to $\Ecu(\rho(a))$.


Let us first observe that by taking an arbitrarily small perturbation of $\rho$ if needed we may assume, in addition to the current hypotheses, that $\Ec(\rho(a))$ is different from $L_0$ and $\Es(\rho(a))$ is not contained in $P_0$. Indeed, there exist arbitrarily small neighborhoods $U$ of $\rho(b)$ so that $U^{(n)}$ is a neighborhood of $\rho(b^{n})$ (Lemma~\ref{lem: power map is open}).
We may take an arbitrarily small perturbation $\rho'(a)$ of $\rho(a)$ so that $\Ec(\rho'(a))$ is different from $L_0$, $\Es(\rho(a))$ does not belong to $P_0$, $L_0$ belongs to $\Ecu(\rho'(a))$, and
\[
\rho'(a)^{-m}\rho(a^{m})\rho(b^{n}) \in U^{(n)}.
\]
In particular, there is some $\rho'(b)\in U$ so that
\[
\rho'(a)^{-m}\rho(a^{m})\rho(b^{n})=\rho'(b)^{n}.
\]
By keeping all other generators of $\Gamma$ fixed, this defines a perturbation $\rho'$ of $\rho$ for which $\rho'(\omega)=\rho(\omega)$, and $\rho'$ is still in the hypotheses of the proposition, with the extra claimed genericity assumptions on $\Ec(\rho'(a))$ and $\Es(\rho'(a))$. To lighten notations, we assume the representation $\rho$ that we started with also satisfies these conditions.



Now fix any continuous path $t\mapsto\rho_t(a)$ passing through $\rho(a)$ at $t=0$, which is constructed only deforming eigenline $\Eu(\rho(a))$ in such a way that $t\mapsto \Eu(\rho_t(a))$ intersects transversely $\Ecu(\rho(a))$ precisely at $t=0$. As above, we may find a continuous path $t\mapsto \rho_t(b)$ passing through $\rho(b)$ at $t=0$ and such that $\rho_t(\omega)=\rho(\omega)$ for all $t$.



Fix a unit vector $v\in L_0$ and let $t\mapsto \theta_t\in(\RR^3)^*\setminus\{0\}$ be a continuous path so that
\[
\ker\theta_t=\Ecu(\rho_t(a))
\]
for all $t$.
As $\Ec(\rho(a))=\Ec(\rho_t(a))\neq L_0$ and $L_0\in\Ecu(\rho(a))$, by our transversality assumption we may find arbitrarily small $t_-<0<t_+$ so that
\[
\theta_{t_-}(v)\cdot \theta_{t_+}(v)<0.
\]
Without loss of generality we may assume
\begin{equation}\label{eq: criterion for non robust qi}
\theta_{t_-}(v)<0<\theta_{t_+}(v).
\end{equation}


On the other hand, for every positive integer $p$ we may take a continuous path $t\mapsto\theta_{t,p}\in(\RR^3)^*\setminus\{0\}$ so that
\[
\ker\theta_{t,p}=\rho_t(a^p)\cdot P_0
\]
holds for all $t$.
Note that, as $\Es(\rho(a))=\Es(\rho_t(a))\not\subset P_0$, for a given $t$ we have
\[
\rho_t(a^p)\cdot P_0\to\Ecu(\rho_t(a)),
\]
as $p\to\infty$.
In particular, we may assume that $\theta_{t,p}$ is chosen in such a way that for every $t$ one has
\[
\theta_{t,p}\to\theta_t
\]
as $p\to\infty$.
By Equation~\eqref{eq: criterion for non robust qi} we may find a large enough $p$ such that
\[
\theta_{t_-,p}(v)<0<\theta_{t_+,p}(v).
\]
There exists then some $t_-<t_0<t_+$ such that $\theta_{t_0,p}(v)=0$ or, in other words,
\[
L_0\subset \rho_{t_0}(a^{p})\cdot P_0.
\]
In particular, $\rho_{t_0}(\omega^q)=\rho(\omega^q)$ preserves the hyperplane $\rho_{t_0}(a^{p})\cdot P_0$, because it preserves $L_0$ and $P_0\cap\rho_{t_0}(a^{p})\cdot P_0$.

On the other hand, the element $\rho_{t_0}(a^p\omega^q a^{-p})$ also preserves $\rho_{t_0}(a^{p})\cdot P_0$ and furthermore acts as a scalar multiple of the identity on it. Hence, the commutator $\rho_{t_0}([\omega^q,a^p\omega^q a^{-p}])$ acts as the identity on it. As $ \rho_{t_0}([\omega^q,a^p\omega^q a^{-p}]) $ has determinant equal to $1$, it is a unipotent element.
\end{proof}

\end{document}
